\scanpage{38}
Soit $\rho\in\Gamma(S,\mathcal O_S)$, $U_\rho\subset\mathbb E^{1}_{S}$ formé des $\lambda$ tels que $1+\lambda\rho$ inversible i.e.

\[
U_\rho=\underline{\operatorname{Spec}}\ \mathcal O_S[T]_{(1+\rho T)}
\]

prenons la multiplication

\[
(\lambda\mathbin{\underset{\rho}{*}}\lambda'=\lambda+\lambda'+\lambda\lambda'\rho
\]

de sorte que $U_\rho$ est un schéma en groupes\inserted{(lisse comm. \unsure{de rel.|de dim.\ rel.} 1)}, et

\[
\begin{aligned}
\lambda&\mapsto 1+\lambda\rho\\
U_\rho&\longrightarrow G_m
\end{aligned}
\]

est un hom de groupes, qui est un iso. sss $\rho$ inversible.

Considérons un hyperpolyèdre régulier \unsure{paramétré|ponctué|pseudo-} épinglé sur $S$, de dim. $n$\inserted{(\unsure{ou 3|ou 1|tous|conv.})}, correspondant \unsure{à|:|;} des invariants $\mu_i$ $(0\leq i\leq n)$ et $\beta_i$ $(0\leq i\leq n-1)$ [avec $\mu_i+1$ inv. $0\leq i\leq n$). Soit \deleted{$\rho=\unsure{\mathbb L|L}_{n+1}(\mu_0,\lambda_0,\ldots,\lambda_{n-1},\mu_n)$}

\[
\unsure{=}\qquad \rho=\mathbb L_{n+1}(\mu_0,\lambda_0,\ldots,\lambda_{n-1},\mu_n)=\langle\pi,\pi'\rangle
\]
% archive sheet number: 37
